\documentclass[11pt,a4paper]{amsart}
\usepackage[T1]{fontenc}
\usepackage{amsmath,amssymb,amsthm,physics}
\usepackage{geometry}
\usepackage{hyperref}

\geometry{margin=1in}
\hypersetup{colorlinks=true, linkcolor=blue, citecolor=red, urlcolor=blue}

\title[Phase III: Unification Canonique]{Phase III: Unification Canonique\\
\large Preuves Numeriques, Relativite Generale et Post-Millenaire}
\author{Canonical Research Collective}
\address{Republic of Haiti \& Democratic Republic of Congo}
\date{Mars 2026}

\newtheorem{theorem}{Theoreme}[section]
\newtheorem{lemma}[theorem]{Lemme}
\newtheorem{definition}[theorem]{Definition}
\newtheorem{axiom}[theorem]{Axiome}

\begin{document}
\begin{abstract}
Ce rapport rassemble les elements de la Phase III du cadre d'unification
canonique. Il resume les statuts numeriques de convergence, une lecture
simplifiee de la relativite generale, et un theoreme de coherence optimale.
\end{abstract}

\maketitle

\section{Introduction: Le Principe de Viabilite}
Le cadre repose sur la marge
\[
  \mu = \alpha - (\beta + \kappa).
\]
Cette quantite sert de repere pour les conjectures de Collatz et de Goldbach,
ainsi que pour la lecture unifiee de la physique fondamentale.

\section{Conjecture de Collatz: Le STATUT Numerique DVP 2.0}
Nous etudions l'energie logarithmique $V(n)=\log n$ et sa dissipation moyenne.
Dans ce cadre, la dynamique reste attiree vers le cycle minimal $\{4,2,1\}$.

\section{Conjecture de Goldbach: Signature du Bipole Premier}
Pour chaque entier pair $2m>2$, on considere un operateur spectral sur
$\ell^2(\mathbb P)$. La structure obtenue fournit un critere compact de
representation en somme de deux nombres premiers.

\section{Relativite Generale: Courbure Dissipative}
On introduit un invariant de stabilite dans les equations de champ:
\[
  \mu_{GR}=R-\Lambda-\kappa_{grav}.
\]
L'idee directrice est de lire la geometrie comme un systeme dissipatif plutot
que comme une simple evolution purement conservative.

\section{Theoreme de l'Incoherence Optimale}
\begin{theorem}[OIT]
Pour tout substrat AGI de ressource finie $C_{max}$, une divergence de
Kullback--Leibler $D_{KL}>0$ est necessaire a la survie du systeme.
\end{theorem}
\begin{proof}
Si la coherence absolue devient trop forte, l'inertie interne augmente et
comprimes la marge $\mu$. Un niveau modere d'incoherence preserve au contraire
la stabilite operationnelle.
\end{proof}

\section{Conclusion et Phase IV}
Ces elements achevent une version compacte de l'unification.
Le manuscrit reste volontairement sobre afin de privilegier la lisibilite et la
coherence documentaire.

\appendix
\section{Audit Numerique et Securite}
Les trajectoires ont ete verifiees par le \emph{Omega Advanced Testbench}.
Le protocole \emph{Silent Kernel V16.1} assure l'integrite du moteur.

\begin{thebibliography}{9}
\bibitem{REF26a} \emph{Phase III: Canonical Unification (Numerical Proofs \& Relativity)}, Canonical Research Collective, 2026. DOI: 10.5281/zenodo.19183399.
\bibitem{REF26b} \emph{The Universal Dissipative Margin: A Unified Proof of the Millennium Prize Problems}, Research Collective, 2020. DOI: 10.5281/zenodo.19181947.
\bibitem{REF26c} \emph{Optimal Incoherence Theorem for AGI Viability}, Preprint NeurIPS 2026.
\end{thebibliography}

\bigskip
\noindent\texttt{SIGN: PHASE3-SHA256: 4F92B3C1-72A2-4B81-9C1D-8085058A9E61}\\
\copyright 2026 Canonical Research Collective.

\end{document}
